
\section*{Histoire : La fracture de 1880 (Ferry/Bréal) et la naissance du dogme philologique}

Cette agonie contemporaine n'est pas un accident. Elle est l'aboutissement logique d'une rupture fondatrice survenue en 1880 : le moment où Jules Ferry et Michel Bréal imposèrent un nouveau paradigme pédagogique pour l'enseignement du grec ancien.

Avant 1880, le grec s'enseignait dans une logique rhétorique, héritière des collèges jésuites. Les élèves écrivaient, composaient en vers grec, imitaient les modèles. Le grec était une langue de production, instrument de formation morale et d'éloquence. C'était une pédagogie de la parole vivante, même si désormais historiquement lointaine.

Après 1880, un changement de régime : le grec devient un objet scientifique, une matière à analyser comme on analyse les structures d'une langue morte au microscope. La « Version » (traduction du grec vers le français) devient l'exercice roi. La grammaire, explicitée, codifiée, systématisée selon les méthodes philologiques allemandes, devient le cœur du curriculum. L'objectif n'est plus de parler ou même de comprendre vivamment, mais d'analyser, de décomposer, de démontrer sa rigueur intellectuelle par la maîtrise de règles abstraites.

Ce changement, porteur de modernité scientifique et d'ambition républicaine, a cristallisé trois dogmes pédagogiques qui dominent encore aujourd'hui : (1) l'idée que le grec « muscle » l'esprit par la gymnastique mentale ; (2) la vénération du Texte-Monument, objet d'analyse à distance ; (3) l'acceptation implicite que seule une élite peut et doit apprendre le grec.
